\documentclass[a4paper,12pt,twoside]{article}
\usepackage[left=1in,right=1in,top=1in,bottom=1in]{geometry}

\usepackage{xeCJK}
\setCJKmainfont{Songti TC}
\setCJKfamilyfont{kai}{Songti TC}
\newcommand{\kai}{\CJKfamily{kai}} % 用 \kai{使用楷書}

\usepackage{lmodern}
\usepackage{titling}
\usepackage{titlesec}
\usepackage{graphicx}
\usepackage{float}
\usepackage{amsmath, amsthm, amssymb, amsfonts, mathrsfs}
\usepackage{upgreek} % 直立的希臘字母

\usepackage[shortlabels]{enumitem}
\usepackage{setspace}
\renewcommand{\baselinestretch}{1.5}
\usepackage[colorlinks]{hyperref}
\usepackage{fancyvrb}
\usepackage{fancyhdr}
\usepackage{multicol}
\usepackage{caption}
\usepackage{listings}
\usepackage{bm}
\usepackage{bbm} % for \mathbbm{1} (indicator function)
\usepackage{centernot}
\usepackage{mdframed}
\usepackage{tikz}
\usetikzlibrary{patterns, arrows.meta, positioning}
\usepackage[nameinlink, capitalise]{cleveref}
\usepackage{diagbox}
\usepackage{booktabs}
\usepackage{multirow,array}
\usepackage{tabularx}
\usepackage{mathtools}
\usepackage{indentfirst} % 首段 indent
\setlength{\parindent}{2em}
\usepackage{cancel} % 數學式子畫線槓掉
\newcommand\mcancel[2][red]{\setbox0=\hbox{$#2$}%
\rlap{\raisebox{.45\ht0}{\textcolor{#1}{\rule{\wd0}{1pt}}}}#2}
\usepackage{datetime2}

\AtBeginDocument{%
 \setlength{\abovedisplayskip}{6pt}
 \setlength{\belowdisplayskip}{6pt}
 \setlength{\abovedisplayshortskip}{4pt}
 \setlength{\belowdisplayshortskip}{4pt}
}

\definecolor{codegreen}{rgb}{0,0.6,0}
\definecolor{codegray}{rgb}{0.5,0.5,0.5}
\definecolor{codepurple}{rgb}{0.58,0,0.82}
\definecolor{backcolour}{rgb}{0.95,0.95,0.92}

\lstdefinestyle{mystyle}{
 backgroundcolor=\color{backcolour},
 commentstyle=\color{codegreen},
 keywordstyle=\color{magenta},
 numberstyle=\tiny\color{codegray},
 stringstyle=\color{codepurple},
 basicstyle=\ttfamily\footnotesize,
 breakatwhitespace=false,
 breaklines=true,
 captionpos=b,
 keepspaces=true,
 numbers=left,
 numbersep=5pt,
 showspaces=false,
 showstringspaces=false,
 showtabs=false,
 tabsize=2
}
\lstset{style=mystyle}

\hypersetup{citecolor=blue, linkcolor=blue, urlcolor=blue}

\DeclareMathOperator{\supp}{supp}
\DeclareMathOperator{\skewness}{Skewness}
\DeclareMathOperator{\DU}{D-U}
\DeclareMathOperator{\CU}{C-U}
\DeclareMathOperator{\Geo}{Geo}
\DeclareMathOperator{\Pois}{Pois}
\DeclareMathOperator{\Gam}{Gam}
\DeclareMathOperator{\LN}{LogN}
\DeclareMathOperator{\cov}{Cov}
\DeclareMathOperator*{\argmax}{arg\,max}
\DeclareMathOperator*{\argmin}{arg\,min}
\DeclareMathOperator{\re}{Re}
\DeclareMathOperator{\im}{Im}

\newcommand{\N}{\mathbb{N}}
\newcommand{\Z}{\mathbb{Z}}
\newcommand{\Q}{\mathbb{Q}}
\newcommand{\R}{\mathbb{R}}
\newcommand{\C}{\mathbb{C}}
\newcommand{\G}{\mathcal{G}}
\newcommand{\mP}{\mathbb{P}}
\newcommand{\F}{\mathcal{F}}
\newcommand{\B}{\mathcal{B}}
\newcommand{\1}{\mathbbm{1}}
\newcommand{\Poi}{\text{Poi}}
\newcommand{\Var}{\text{Var}}
\newcommand{\as}{\text{a.s.}}
\newcommand{\vect}[1]{\boldsymbol{\mathbf{#1}}}
\newcommand{\paren}[1]{\left({#1}\right)}
\newcommand{\brk}[1]{\left[{#1}\right]}
\newcommand{\curl}[1]{\left\{{#1}\right\}}
\renewcommand{\lambda}{\lambdaup}

\theoremstyle{definition} % 定理環境 style 改用正體
\newtheorem{theorem}{定理}
\newtheorem{definition}{定義}
\newtheorem*{definition*}{定義}
\newtheorem{example}{例子}
\newtheorem{exercise}{練習}
\newtheorem{corollary}{推論}
\newtheorem{property}{性質}
\newtheorem{prop}{命題}
\newtheorem{lemma}{引理}
\newtheorem{problem}{問題}
\newtheorem*{remark}{註}

\renewcommand{\figurename}{圖}
\renewcommand{\tablename}{表}

\DeclareMathOperator{\var}{Var}
\DeclareMathOperator{\corr}{corr}
\DeclareMathOperator{\rank}{rank}
\DeclareMathOperator{\MSE}{MSE}
\DeclareMathOperator{\SSE}{SSE}
\DeclareMathOperator{\SST}{SST}
\DeclareMathOperator{\SSR}{SSR}
\DeclareMathOperator{\sgn}{sign}
\DeclareMathOperator{\alt}{Alt}
\DeclareMathOperator{\adj}{adj}
\DeclareMathOperator{\curlOp}{curl}
\DeclareMathOperator{\spn}{span}
\DeclareMathOperator{\sse}{SSE}
\DeclareMathOperator{\tr}{tr}

\newcommand{\ep}{\varepsilon}
\newcommand{\p}{\partial}
\newcommand{\abs}[1]{\left|#1\right|}
\newcommand\norm[1]{\left\lVert#1\right\rVert}
\newcommand{\innerprod}[1]{\langle#1\rangle}
\newcommand{\ds}{\displaystyle}
\newcommand{\blank}{\rule[-1mm]{1.5cm}{0.1mm}}
\newcommand{\dint}{\displaystyle\int}
\renewcommand{\bar}{\overline}

\newtheoremstyle{remark}{}{}{}{}{\bfseries\sffamily}{.}{.5em}{#1}
\theoremstyle{remark}

\newtheorem*{remark*}{\textit{Remark}}
\newtheorem*{solution*}{解}

\AtEndEnvironment{solution*}{\hfill$\square$}
\newtheorem*{hint*}{Hint}


\newcommand{\redtext}[1]{{\color{red} #1}}
\newcommand{\bluetext}[1]{{\color{blue} #1}}

\fancypagestyle{firstpage}{
  \fancyhf{}
  \fancyhead[L]{\small 《高中數學實驗教材》第四冊（自然組）}
  \fancyhead[R]{\small 附錄四\quad 對數表}
  \fancyfoot[C]{\thepage}
}

\pagestyle{fancy}
\setlength{\headheight}{14pt}
\renewcommand{\headrulewidth}{0.5pt}
\renewcommand{\footrulewidth}{0pt}
\fancyhead{}
\fancyhead[LE]{\small 《高中數學實驗教材》第四冊（自然組）}
\fancyhead[RO]{\small 附錄四\quad 對數表}
\fancyfoot[C]{\thepage}

\title{\LARGE\textbf{《高中數學實驗教材》第四冊（自然組）}\\[0.5em] \Large\textbf{附錄四\quad 對數表}}
\author{\textbf{高中數學實驗教材編輯小組}\quad 編著}
\date{民國六十五年（1976）三版}

\begin{document}

\maketitle
\thispagestyle{firstpage}

\clearpage
\tableofcontents
\newpage

\section*{附錄四\quad 對數表}
\addcontentsline{toc}{section}{附錄四\texorpdfstring{\quad}{ }對數表}
\label{sec:b4_a4}

\subsection*{表 1.\quad 常用對數表}
\addcontentsline{toc}{subsection}{表 1.\texorpdfstring{\quad}{ }常用對數表}
\begin{center} $y = \log_{10} x$ \end{center}
{\centering\fontsize{7.5pt}{8.8pt}\selectfont\setlength{\tabcolsep}{1.5pt}\renewcommand{\arraystretch}{0.96}
\begin{tabular}{|c|cccccccccc|ccccccccc|}
\hline
x & 0 & 1 & 2 & 3 & 4 & 5 & 6 & 7 & 8 & 9 & 1 & 2 & 3 & 4 & 5 & 6 & 7 & 8 & 9 \\ \hline
1.0 & .0000 & .0043 & .0086 & .0128 & .0170 & .0212 & .0253 & .0294 & .0334 & .0374 & 4 & 8 & 12 & 17 & 21 & 25 & 29 & 33 & 37 \\
1.1 & .0414 & .0453 & .0492 & .0531 & .0569 & .0607 & .0645 & .0682 & .0719 & .0755 & 4 & 8 & 11 & 15 & 19 & 23 & 26 & 30 & 34 \\
1.2 & .0792 & .0828 & .0864 & .0899 & .0934 & .0969 & .1004 & .1038 & .1072 & .1106 & 3 & 7 & 10 & 14 & 17 & 21 & 24 & 28 & 31 \\
1.3 & .1139 & .1173 & .1206 & .1239 & .1271 & .1303 & .1335 & .1367 & .1399 & .1430 & 3 & 6 & 10 & 13 & 16 & 19 & 23 & 26 & 29 \\
1.4 & .1461 & .1492 & .1523 & .1553 & .1584 & .1614 & .1644 & .1673 & .1703 & .1732 & 3 & 6 & 9 & 12 & 15 & 18 & 21 & 24 & 27 \\
1.5 & .1761 & .1790 & .1818 & .1847 & .1875 & .1903 & .1931 & .1959 & .1987 & .2014 & 3 & 6 & 8 & 11 & 14 & 17 & 20 & 22 & 25 \\
1.6 & .2041 & .2068 & .2095 & .2122 & .2148 & .2175 & .2201 & .2227 & .2253 & .2279 & 3 & 5 & 8 & 11 & 13 & 16 & 18 & 21 & 24 \\
1.7 & .2304 & .2330 & .2355 & .2380 & .2405 & .2430 & .2455 & .2480 & .2504 & .2529 & 2 & 5 & 7 & 10 & 12 & 15 & 17 & 20 & 22 \\
1.8 & .2553 & .2577 & .2601 & .2625 & .2648 & .2672 & .2695 & .2718 & .2742 & .2765 & 2 & 5 & 7 & 9 & 12 & 14 & 16 & 19 & 21 \\
1.9 & .2788 & .2810 & .2833 & .2856 & .2878 & .2900 & .2923 & .2945 & .2967 & .2989 & 2 & 4 & 7 & 9 & 11 & 13 & 16 & 18 & 20 \\
2.0 & .3010 & .3032 & .3054 & .3075 & .3096 & .3118 & .3139 & .3160 & .3181 & .3201 & 2 & 4 & 6 & 8 & 11 & 13 & 15 & 17 & 19 \\
2.1 & .3222 & .3243 & .3263 & .3284 & .3304 & .3324 & .3345 & .3365 & .3385 & .3404 & 2 & 4 & 6 & 8 & 10 & 12 & 14 & 16 & 18 \\
2.2 & .3424 & .3444 & .3464 & .3483 & .3502 & .3522 & .3541 & .3560 & .3579 & .3598 & 2 & 4 & 6 & 8 & 10 & 12 & 14 & 15 & 17 \\
2.3 & .3617 & .3636 & .3655 & .3674 & .3692 & .3711 & .3729 & .3747 & .3766 & .3784 & 2 & 4 & 6 & 7 & 9 & 11 & 13 & 15 & 17 \\
2.4 & .3802 & .3820 & .3838 & .3856 & .3874 & .3892 & .3909 & .3927 & .3945 & .3962 & 2 & 4 & 5 & 7 & 9 & 11 & 12 & 14 & 16 \\
2.5 & .3979 & .3997 & .4014 & .4031 & .4048 & .4065 & .4082 & .4099 & .4116 & .4133 & 2 & 3 & 5 & 7 & 9 & 10 & 12 & 14 & 15 \\
2.6 & .4150 & .4166 & .4183 & .4200 & .4216 & .4232 & .4249 & .4265 & .4281 & .4298 & 2 & 3 & 5 & 7 & 8 & 10 & 11 & 13 & 15 \\
2.7 & .4314 & .4330 & .4346 & .4362 & .4378 & .4393 & .4409 & .4425 & .4440 & .4456 & 2 & 3 & 5 & 6 & 8 & 9 & 11 & 13 & 14 \\
2.8 & .4472 & .4487 & .4502 & .4518 & .4533 & .4548 & .4564 & .4579 & .4594 & .4609 & 2 & 3 & 5 & 6 & 8 & 9 & 11 & 12 & 14 \\
2.9 & .4624 & .4639 & .4654 & .4669 & .4683 & .4698 & .4713 & .4728 & .4742 & .4757 & 1 & 3 & 4 & 6 & 7 & 9 & 10 & 12 & 13 \\
3.0 & .4771 & .4786 & .4800 & .4814 & .4829 & .4843 & .4857 & .4871 & .4886 & .4900 & 1 & 3 & 4 & 6 & 7 & 9 & 10 & 11 & 13 \\
3.1 & .4914 & .4928 & .4942 & .4955 & .4969 & .4983 & .4997 & .5011 & .5024 & .5038 & 1 & 3 & 4 & 6 & 7 & 8 & 10 & 11 & 12 \\
3.2 & .5051 & .5065 & .5079 & .5092 & .5105 & .5119 & .5132 & .5145 & .5159 & .5172 & 1 & 3 & 4 & 5 & 7 & 8 & 9 & 11 & 12 \\
3.3 & .5185 & .5198 & .5211 & .5224 & .5237 & .5250 & .5263 & .5276 & .5289 & .5302 & 1 & 3 & 4 & 5 & 6 & 8 & 9 & 10 & 12 \\
3.4 & .5315 & .5328 & .5340 & .5353 & .5366 & .5378 & .5391 & .5403 & .5416 & .5428 & 1 & 3 & 4 & 5 & 6 & 8 & 9 & 10 & 11 \\
3.5 & .5441 & .5453 & .5465 & .5478 & .5490 & .5502 & .5514 & .5527 & .5539 & .5551 & 1 & 2 & 4 & 5 & 6 & 7 & 9 & 10 & 11 \\
3.6 & .5563 & .5575 & .5587 & .5599 & .5611 & .5623 & .5635 & .5647 & .5658 & .5670 & 1 & 2 & 4 & 5 & 6 & 7 & 8 & 10 & 11 \\
3.7 & .5682 & .5694 & .5705 & .5717 & .5729 & .5740 & .5752 & .5763 & .5775 & .5786 & 1 & 2 & 3 & 5 & 6 & 7 & 8 & 9 & 10 \\
3.8 & .5798 & .5809 & .5821 & .5832 & .5843 & .5855 & .5866 & .5877 & .5888 & .5899 & 1 & 2 & 3 & 5 & 6 & 7 & 8 & 9 & 10 \\
3.9 & .5911 & .5922 & .5933 & .5944 & .5955 & .5966 & .5977 & .5988 & .5999 & .6010 & 1 & 2 & 3 & 4 & 5 & 7 & 8 & 9 & 10 \\
4.0 & .6021 & .6031 & .6042 & .6053 & .6064 & .6075 & .6085 & .6096 & .6107 & .6117 & 1 & 2 & 3 & 4 & 5 & 6 & 8 & 9 & 10 \\
4.1 & .6128 & .6138 & .6149 & .6160 & .6170 & .6180 & .6191 & .6201 & .6212 & .6222 & 1 & 2 & 3 & 4 & 5 & 6 & 7 & 8 & 9 \\
4.2 & .6232 & .6243 & .6253 & .6263 & .6274 & .6284 & .6294 & .6304 & .6314 & .6325 & 1 & 2 & 3 & 4 & 5 & 6 & 7 & 8 & 9 \\
4.3 & .6335 & .6345 & .6355 & .6365 & .6375 & .6385 & .6395 & .6405 & .6415 & .6425 & 1 & 2 & 3 & 4 & 5 & 6 & 7 & 8 & 9 \\
4.4 & .6435 & .6444 & .6454 & .6464 & .6474 & .6484 & .6493 & .6503 & .6513 & .6522 & 1 & 2 & 3 & 4 & 5 & 6 & 7 & 8 & 9 \\
4.5 & .6532 & .6542 & .6551 & .6561 & .6571 & .6580 & .6590 & .6599 & .6609 & .6618 & 1 & 2 & 3 & 4 & 5 & 6 & 7 & 8 & 9 \\
4.6 & .6628 & .6637 & .6646 & .6656 & .6665 & .6675 & .6684 & .6693 & .6702 & .6712 & 1 & 2 & 3 & 4 & 5 & 6 & 7 & 7 & 8 \\
4.7 & .6721 & .6730 & .6739 & .6749 & .6758 & .6767 & .6776 & .6785 & .6794 & .6803 & 1 & 2 & 3 & 4 & 5 & 5 & 6 & 7 & 8 \\
4.8 & .6812 & .6821 & .6830 & .6839 & .6848 & .6857 & .6866 & .6875 & .6884 & .6893 & 1 & 2 & 3 & 4 & 4 & 5 & 6 & 7 & 8 \\
4.9 & .6902 & .6911 & .6920 & .6928 & .6937 & .6946 & .6955 & .6964 & .6972 & .6981 & 1 & 2 & 3 & 4 & 4 & 5 & 6 & 7 & 8 \\
5.0 & .6990 & .6998 & .7007 & .7016 & .7024 & .7033 & .7042 & .7050 & .7059 & .7067 & 1 & 2 & 3 & 3 & 4 & 5 & 6 & 7 & 8 \\
5.1 & .7076 & .7084 & .7093 & .7101 & .7110 & .7118 & .7126 & .7135 & .7143 & .7152 & 1 & 2 & 3 & 3 & 4 & 5 & 6 & 7 & 8 \\
5.2 & .7160 & .7168 & .7177 & .7185 & .7193 & .7202 & .7210 & .7218 & .7226 & .7235 & 1 & 2 & 2 & 3 & 4 & 5 & 6 & 7 & 7 \\
5.3 & .7243 & .7251 & .7259 & .7267 & .7275 & .7284 & .7292 & .7300 & .7308 & .7316 & 1 & 2 & 2 & 3 & 4 & 5 & 6 & 6 & 7 \\
5.4 & .7324 & .7332 & .7340 & .7348 & .7356 & .7364 & .7372 & .7380 & .7388 & .7396 & 1 & 2 & 2 & 3 & 4 & 5 & 6 & 6 & 7 \\
\hline
\end{tabular}
\par}
\clearpage

\subsection*{表 1.\quad （續）常用對數表}
\begin{center} $y = \log_{10} x$ \end{center}
{\centering\fontsize{7.5pt}{8.8pt}\selectfont\setlength{\tabcolsep}{1.5pt}\renewcommand{\arraystretch}{0.96}
\begin{tabular}{|c|cccccccccc|ccccccccc|}
\hline
x & 0 & 1 & 2 & 3 & 4 & 5 & 6 & 7 & 8 & 9 & 1 & 2 & 3 & 4 & 5 & 6 & 7 & 8 & 9 \\ \hline
5.5 & .7404 & .7412 & .7419 & .7427 & .7435 & .7443 & .7451 & .7459 & .7466 & .7474 & 1 & 2 & 2 & 3 & 4 & 5 & 5 & 6 & 7 \\
5.6 & .7482 & .7490 & .7497 & .7505 & .7513 & .7520 & .7528 & .7536 & .7543 & .7551 & 1 & 2 & 2 & 3 & 4 & 5 & 5 & 6 & 7 \\
5.7 & .7559 & .7566 & .7574 & .7582 & .7589 & .7597 & .7604 & .7612 & .7619 & .7627 & 1 & 2 & 2 & 3 & 4 & 5 & 5 & 6 & 7 \\
5.8 & .7634 & .7642 & .7649 & .7657 & .7664 & .7672 & .7679 & .7686 & .7694 & .7701 & 1 & 1 & 2 & 3 & 4 & 4 & 5 & 6 & 7 \\
5.9 & .7709 & .7716 & .7723 & .7731 & .7738 & .7745 & .7752 & .7760 & .7767 & .7774 & 1 & 1 & 2 & 3 & 4 & 4 & 5 & 6 & 7 \\
6.0 & .7782 & .7789 & .7796 & .7803 & .7810 & .7818 & .7825 & .7832 & .7839 & .7846 & 1 & 1 & 2 & 3 & 4 & 4 & 5 & 6 & 6 \\
6.1 & .7853 & .7860 & .7868 & .7875 & .7882 & .7889 & .7896 & .7903 & .7910 & .7917 & 1 & 1 & 2 & 3 & 4 & 4 & 5 & 6 & 6 \\
6.2 & .7924 & .7931 & .7938 & .7945 & .7952 & .7959 & .7966 & .7973 & .7980 & .7987 & 1 & 1 & 2 & 3 & 3 & 4 & 5 & 6 & 6 \\
6.3 & .7993 & .8000 & .8007 & .8014 & .8021 & .8028 & .8035 & .8041 & .8048 & .8055 & 1 & 1 & 2 & 3 & 3 & 4 & 5 & 5 & 6 \\
6.4 & .8062 & .8069 & .8075 & .8082 & .8089 & .8096 & .8102 & .8109 & .8116 & .8122 & 1 & 1 & 2 & 3 & 3 & 4 & 5 & 5 & 6 \\
6.5 & .8129 & .8136 & .8142 & .8149 & .8156 & .8162 & .8169 & .8176 & .8182 & .8189 & 1 & 1 & 2 & 3 & 3 & 4 & 5 & 5 & 6 \\
6.6 & .8195 & .8202 & .8209 & .8215 & .8222 & .8228 & .8235 & .8241 & .8248 & .8254 & 1 & 1 & 2 & 3 & 3 & 4 & 5 & 5 & 6 \\
6.7 & .8261 & .8267 & .8274 & .8280 & .8287 & .8293 & .8299 & .8306 & .8312 & .8319 & 1 & 1 & 2 & 3 & 3 & 4 & 5 & 5 & 6 \\
6.8 & .8325 & .8331 & .8338 & .8344 & .8351 & .8357 & .8363 & .8370 & .8376 & .8382 & 1 & 1 & 2 & 3 & 3 & 4 & 4 & 5 & 6 \\
6.9 & .8388 & .8395 & .8401 & .8407 & .8414 & .8420 & .8426 & .8432 & .8439 & .8445 & 1 & 1 & 2 & 2 & 3 & 4 & 4 & 5 & 6 \\
7.0 & .8451 & .8457 & .8463 & .8470 & .8476 & .8482 & .8488 & .8494 & .8500 & .8506 & 1 & 1 & 2 & 2 & 3 & 4 & 4 & 5 & 6 \\
7.1 & .8513 & .8519 & .8525 & .8531 & .8537 & .8543 & .8549 & .8555 & .8561 & .8567 & 1 & 1 & 2 & 2 & 3 & 4 & 4 & 5 & 5 \\
7.2 & .8573 & .8579 & .8585 & .8591 & .8597 & .8603 & .8609 & .8615 & .8621 & .8627 & 1 & 1 & 2 & 2 & 3 & 4 & 4 & 5 & 5 \\
7.3 & .8633 & .8639 & .8645 & .8651 & .8657 & .8663 & .8669 & .8675 & .8681 & .8686 & 1 & 1 & 2 & 2 & 3 & 4 & 4 & 5 & 5 \\
7.4 & .8692 & .8698 & .8704 & .8710 & .8716 & .8722 & .8727 & .8733 & .8739 & .8745 & 1 & 1 & 2 & 2 & 3 & 3 & 4 & 5 & 5 \\
7.5 & .8751 & .8756 & .8762 & .8768 & .8774 & .8779 & .8785 & .8791 & .8797 & .8802 & 1 & 1 & 2 & 2 & 3 & 3 & 4 & 5 & 5 \\
7.6 & .8808 & .8814 & .8820 & .8825 & .8831 & .8837 & .8842 & .8848 & .8854 & .8859 & 1 & 1 & 2 & 2 & 3 & 3 & 4 & 5 & 5 \\
7.7 & .8865 & .8871 & .8876 & .8882 & .8887 & .8893 & .8899 & .8904 & .8910 & .8915 & 1 & 1 & 2 & 2 & 3 & 3 & 4 & 4 & 5 \\
7.8 & .8921 & .8927 & .8932 & .8938 & .8943 & .8949 & .8954 & .8960 & .8965 & .8971 & 1 & 1 & 2 & 2 & 3 & 3 & 4 & 4 & 5 \\
7.9 & .8976 & .8982 & .8987 & .8993 & .8998 & .9004 & .9009 & .9015 & .9020 & .9025 & 1 & 1 & 2 & 2 & 3 & 3 & 4 & 4 & 5 \\
8.0 & .9031 & .9036 & .9042 & .9047 & .9053 & .9058 & .9063 & .9069 & .9074 & .9079 & 1 & 1 & 2 & 2 & 3 & 3 & 4 & 4 & 5 \\
8.1 & .9085 & .9090 & .9096 & .9101 & .9106 & .9112 & .9117 & .9122 & .9128 & .9133 & 1 & 1 & 2 & 2 & 3 & 3 & 4 & 4 & 5 \\
8.2 & .9138 & .9143 & .9149 & .9154 & .9159 & .9165 & .9170 & .9175 & .9180 & .9186 & 1 & 1 & 2 & 2 & 3 & 3 & 4 & 4 & 5 \\
8.3 & .9191 & .9196 & .9201 & .9206 & .9212 & .9217 & .9222 & .9227 & .9232 & .9238 & 1 & 1 & 2 & 2 & 3 & 3 & 4 & 4 & 5 \\
8.4 & .9243 & .9248 & .9253 & .9258 & .9263 & .9269 & .9274 & .9279 & .9284 & .9289 & 1 & 1 & 2 & 2 & 3 & 3 & 4 & 4 & 5 \\
8.5 & .9294 & .9299 & .9304 & .9309 & .9315 & .9320 & .9325 & .9330 & .9335 & .9340 & 1 & 1 & 2 & 2 & 3 & 3 & 4 & 4 & 5 \\
8.6 & .9345 & .9350 & .9355 & .9360 & .9365 & .9370 & .9375 & .9380 & .9385 & .9390 & 1 & 1 & 2 & 2 & 3 & 3 & 4 & 4 & 5 \\
8.7 & .9395 & .9400 & .9405 & .9410 & .9415 & .9420 & .9425 & .9430 & .9435 & .9440 & 0 & 1 & 1 & 2 & 2 & 3 & 3 & 4 & 4 \\
8.8 & .9445 & .9450 & .9455 & .9460 & .9465 & .9469 & .9474 & .9479 & .9484 & .9489 & 0 & 1 & 1 & 2 & 2 & 3 & 3 & 4 & 4 \\
8.9 & .9494 & .9499 & .9504 & .9509 & .9513 & .9518 & .9523 & .9528 & .9533 & .9538 & 0 & 1 & 1 & 2 & 2 & 3 & 3 & 4 & 4 \\
9.0 & .9542 & .9547 & .9552 & .9557 & .9562 & .9566 & .9571 & .9576 & .9581 & .9586 & 0 & 1 & 1 & 2 & 2 & 3 & 3 & 4 & 4 \\
9.1 & .9590 & .9595 & .9600 & .9605 & .9609 & .9614 & .9619 & .9624 & .9628 & .9633 & 0 & 1 & 1 & 2 & 2 & 3 & 3 & 4 & 4 \\
9.2 & .9638 & .9643 & .9647 & .9652 & .9657 & .9661 & .9666 & .9671 & .9675 & .9680 & 0 & 1 & 1 & 2 & 2 & 3 & 3 & 4 & 4 \\
9.3 & .9685 & .9689 & .9694 & .9699 & .9703 & .9708 & .9713 & .9717 & .9722 & .9727 & 0 & 1 & 1 & 2 & 2 & 3 & 3 & 4 & 4 \\
9.4 & .9731 & .9736 & .9741 & .9745 & .9750 & .9754 & .9759 & .9763 & .9768 & .9773 & 0 & 1 & 1 & 2 & 2 & 3 & 3 & 4 & 4 \\
9.5 & .9777 & .9782 & .9786 & .9791 & .9795 & .9800 & .9805 & .9809 & .9814 & .9818 & 0 & 1 & 1 & 2 & 2 & 3 & 3 & 4 & 4 \\
9.6 & .9823 & .9827 & .9832 & .9836 & .9841 & .9845 & .9850 & .9854 & .9859 & .9863 & 0 & 1 & 1 & 2 & 2 & 3 & 3 & 4 & 4 \\
9.7 & .9868 & .9872 & .9877 & .9881 & .9886 & .9890 & .9894 & .9899 & .9903 & .9908 & 0 & 1 & 1 & 2 & 2 & 3 & 3 & 4 & 4 \\
9.8 & .9912 & .9917 & .9921 & .9926 & .9930 & .9934 & .9939 & .9943 & .9948 & .9952 & 0 & 1 & 1 & 2 & 2 & 3 & 3 & 4 & 4 \\
9.9 & .9956 & .9961 & .9965 & .9969 & .9974 & .9978 & .9983 & .9987 & .9991 & .9996 & 0 & 1 & 1 & 2 & 2 & 3 & 3 & 3 & 4 \\
\hline
\end{tabular}
\par}
\clearpage

\subsection*{表 2.\quad 自然對數表}
\addcontentsline{toc}{subsection}{表 2.\texorpdfstring{\quad}{ }自然對數表}
\begin{center} $y = \ln x$ \end{center}
{\centering\fontsize{7.5pt}{8.8pt}\selectfont\setlength{\tabcolsep}{3.5pt}\renewcommand{\arraystretch}{0.96}
\begin{tabular}{|c|cccccccccc|}
\hline
$x$ & 0 & 0.01 & 0.02 & 0.03 & 0.04 & 0.05 & 0.06 & 0.07 & 0.08 & 0.09 \\ \hline
1.0 & 0.0000 & 0.0100 & 0.0198 & 0.0296 & 0.0392 & 0.0488 & 0.0583 & 0.0677 & 0.0770 & 0.0862 \\
1.1 & 0.0953 & 0.1044 & 0.1133 & 0.1222 & 0.1310 & 0.1398 & 0.1484 & 0.1570 & 0.1655 & 0.1740 \\
1.2 & 0.1823 & 0.1906 & 0.1989 & 0.2070 & 0.2151 & 0.2231 & 0.2311 & 0.2390 & 0.2469 & 0.2546 \\
1.3 & 0.2624 & 0.2700 & 0.2776 & 0.2852 & 0.2927 & 0.3001 & 0.3075 & 0.3148 & 0.3221 & 0.3293 \\
1.4 & 0.3365 & 0.3436 & 0.3507 & 0.3577 & 0.3646 & 0.3716 & 0.3784 & 0.3853 & 0.3920 & 0.3988 \\
\hline
1.5 & 0.4055 & 0.4121 & 0.4187 & 0.4253 & 0.4318 & 0.4383 & 0.4447 & 0.4511 & 0.4574 & 0.4637 \\
1.6 & 0.4700 & 0.4762 & 0.4824 & 0.4886 & 0.4947 & 0.5008 & 0.5068 & 0.5128 & 0.5188 & 0.5247 \\
1.7 & 0.5306 & 0.5365 & 0.5423 & 0.5481 & 0.5539 & 0.5596 & 0.5653 & 0.5710 & 0.5766 & 0.5822 \\
1.8 & 0.5878 & 0.5933 & 0.5988 & 0.6043 & 0.6098 & 0.6152 & 0.6206 & 0.6259 & 0.6313 & 0.6366 \\
1.9 & 0.6419 & 0.6471 & 0.6523 & 0.6575 & 0.6627 & 0.6678 & 0.6729 & 0.6780 & 0.6831 & 0.6881 \\
\hline
2.0 & 0.6931 & 0.6981 & 0.7031 & 0.7080 & 0.7129 & 0.7178 & 0.7227 & 0.7275 & 0.7324 & 0.7372 \\
2.1 & 0.7419 & 0.7467 & 0.7514 & 0.7561 & 0.7608 & 0.7655 & 0.7701 & 0.7747 & 0.7793 & 0.7839 \\
2.2 & 0.7885 & 0.7930 & 0.7975 & 0.8020 & 0.8065 & 0.8109 & 0.8154 & 0.8198 & 0.8242 & 0.8286 \\
2.3 & 0.8329 & 0.8372 & 0.8416 & 0.8459 & 0.8502 & 0.8544 & 0.8587 & 0.8629 & 0.8671 & 0.8713 \\
2.4 & 0.8755 & 0.8796 & 0.8838 & 0.8879 & 0.8920 & 0.8961 & 0.9002 & 0.9042 & 0.9083 & 0.9123 \\
\hline
2.5 & 0.9163 & 0.9203 & 0.9243 & 0.9282 & 0.9322 & 0.9361 & 0.9400 & 0.9439 & 0.9478 & 0.9517 \\
2.6 & 0.9555 & 0.9594 & 0.9632 & 0.9670 & 0.9708 & 0.9746 & 0.9783 & 0.9821 & 0.9858 & 0.9895 \\
2.7 & 0.9933 & 0.9969 & 1.0006 & 1.0043 & 1.0080 & 1.0116 & 1.0152 & 1.0188 & 1.0225 & 1.0260 \\
2.8 & 1.0296 & 1.0332 & 1.0367 & 1.0403 & 1.0438 & 1.0473 & 1.0508 & 1.0543 & 1.0578 & 1.0613 \\
2.9 & 1.0647 & 1.0682 & 1.0716 & 1.0750 & 1.0784 & 1.0818 & 1.0852 & 1.0886 & 1.0919 & 1.0953 \\
\hline
3.0 & 1.0986 & 1.1019 & 1.1053 & 1.1086 & 1.1119 & 1.1151 & 1.1184 & 1.1217 & 1.1249 & 1.1282 \\
3.1 & 1.1314 & 1.1346 & 1.1378 & 1.1410 & 1.1442 & 1.1474 & 1.1506 & 1.1537 & 1.1569 & 1.1600 \\
3.2 & 1.1632 & 1.1663 & 1.1694 & 1.1725 & 1.1756 & 1.1787 & 1.1817 & 1.1848 & 1.1878 & 1.1909 \\
3.3 & 1.1939 & 1.1969 & 1.2000 & 1.2030 & 1.2060 & 1.2090 & 1.2119 & 1.2149 & 1.2179 & 1.2208 \\
3.4 & 1.2238 & 1.2267 & 1.2296 & 1.2326 & 1.2355 & 1.2384 & 1.2413 & 1.2442 & 1.2470 & 1.2499 \\
\hline
3.5 & 1.2528 & 1.2556 & 1.2585 & 1.2613 & 1.2641 & 1.2669 & 1.2698 & 1.2726 & 1.2754 & 1.2782 \\
3.6 & 1.2809 & 1.2837 & 1.2865 & 1.2892 & 1.2920 & 1.2947 & 1.2975 & 1.3002 & 1.3029 & 1.3056 \\
3.7 & 1.3083 & 1.3110 & 1.3137 & 1.3164 & 1.3191 & 1.3218 & 1.3244 & 1.3271 & 1.3297 & 1.3324 \\
3.8 & 1.3350 & 1.3376 & 1.3403 & 1.3429 & 1.3455 & 1.3481 & 1.3507 & 1.3533 & 1.3558 & 1.3584 \\
3.9 & 1.3610 & 1.3635 & 1.3661 & 1.3686 & 1.3712 & 1.3737 & 1.3762 & 1.3788 & 1.3813 & 1.3838 \\
\hline
4.0 & 1.3863 & 1.3888 & 1.3913 & 1.3938 & 1.3962 & 1.3987 & 1.4012 & 1.4036 & 1.4061 & 1.4085 \\
4.1 & 1.4110 & 1.4134 & 1.4159 & 1.4183 & 1.4207 & 1.4231 & 1.4255 & 1.4279 & 1.4303 & 1.4327 \\
4.2 & 1.4351 & 1.4375 & 1.4398 & 1.4422 & 1.4446 & 1.4469 & 1.4493 & 1.4516 & 1.4540 & 1.4563 \\
4.3 & 1.4586 & 1.4609 & 1.4633 & 1.4656 & 1.4679 & 1.4702 & 1.4725 & 1.4748 & 1.4770 & 1.4793 \\
4.4 & 1.4816 & 1.4839 & 1.4861 & 1.4884 & 1.4907 & 1.4929 & 1.4951 & 1.4974 & 1.4996 & 1.5019 \\
\hline
4.5 & 1.5041 & 1.5063 & 1.5085 & 1.5107 & 1.5129 & 1.5151 & 1.5173 & 1.5195 & 1.5217 & 1.5239 \\
4.6 & 1.5261 & 1.5282 & 1.5304 & 1.5326 & 1.5347 & 1.5369 & 1.5390 & 1.5412 & 1.5433 & 1.5454 \\
4.7 & 1.5476 & 1.5497 & 1.5518 & 1.5539 & 1.5560 & 1.5581 & 1.5602 & 1.5623 & 1.5644 & 1.5665 \\
4.8 & 1.5686 & 1.5707 & 1.5728 & 1.5748 & 1.5769 & 1.5790 & 1.5810 & 1.5831 & 1.5851 & 1.5872 \\
4.9 & 1.5892 & 1.5913 & 1.5933 & 1.5953 & 1.5974 & 1.5994 & 1.6014 & 1.6034 & 1.6054 & 1.6074 \\
\hline
5.0 & 1.6094 & 1.6114 & 1.6134 & 1.6154 & 1.6174 & 1.6194 & 1.6214 & 1.6233 & 1.6253 & 1.6273 \\
5.1 & 1.6292 & 1.6312 & 1.6332 & 1.6351 & 1.6371 & 1.6390 & 1.6409 & 1.6429 & 1.6448 & 1.6467 \\
5.2 & 1.6487 & 1.6506 & 1.6525 & 1.6544 & 1.6563 & 1.6582 & 1.6601 & 1.6620 & 1.6639 & 1.6658 \\
5.3 & 1.6677 & 1.6696 & 1.6715 & 1.6734 & 1.6752 & 1.6771 & 1.6790 & 1.6808 & 1.6827 & 1.6845 \\
5.4 & 1.6864 & 1.6882 & 1.6901 & 1.6919 & 1.6938 & 1.6956 & 1.6974 & 1.6993 & 1.7011 & 1.7029 \\
\hline
$x$ & 0 & 0.01 & 0.02 & 0.03 & 0.04 & 0.05 & 0.06 & 0.07 & 0.08 & 0.09 \\ \hline
\end{tabular}\par}

\vspace{1em}
{\centering\fontsize{8pt}{9.5pt}\selectfont\setlength{\tabcolsep}{4.5pt}\renewcommand{\arraystretch}{1.1}
\begin{tabular}{|c|ccccccccc|}
\hline
$n$ & 1 & 2 & 3 & 4 & 5 & 6 & 7 & 8 & 9 \\ \hline
$\ln 10^n$ & 2.3026 & 4.6052 & 6.9078 & 9.2103 & 11.5129 & 13.8155 & 16.1181 & 18.4207 & 20.7233 \\
$\ln 10^{-n}$ & $\bar{3}.6974$ & $\bar{5}.3948$ & $\bar{7}.0922$ & $\bar{10}.7897$ & $\bar{12}.4874$ & $\bar{14}.1845$ & $\bar{17}.8819$ & $\bar{19}.5793$ & $\bar{21}.2767$ \\ \hline
\end{tabular}\par}

\clearpage

\subsection*{表 2.\quad （續）自然對數表}
\begin{center} $y = \ln x$ \end{center}
{\centering\fontsize{7.5pt}{8.8pt}\selectfont\setlength{\tabcolsep}{3.5pt}\renewcommand{\arraystretch}{0.96}
\begin{tabular}{|c|cccccccccc|}
\hline
$x$ & 0 & 0.01 & 0.02 & 0.03 & 0.04 & 0.05 & 0.06 & 0.07 & 0.08 & 0.09 \\ \hline
5.5 & 1.7047 & 1.7066 & 1.7084 & 1.7102 & 1.7120 & 1.7138 & 1.7156 & 1.7174 & 1.7192 & 1.7210 \\
5.6 & 1.7228 & 1.7246 & 1.7263 & 1.7281 & 1.7299 & 1.7317 & 1.7334 & 1.7352 & 1.7370 & 1.7387 \\
5.7 & 1.7405 & 1.7422 & 1.7440 & 1.7457 & 1.7475 & 1.7492 & 1.7509 & 1.7527 & 1.7544 & 1.7561 \\
5.8 & 1.7579 & 1.7596 & 1.7613 & 1.7630 & 1.7647 & 1.7664 & 1.7681 & 1.7699 & 1.7716 & 1.7733 \\
5.9 & 1.7750 & 1.7766 & 1.7783 & 1.7800 & 1.7817 & 1.7834 & 1.7851 & 1.7867 & 1.7884 & 1.7901 \\
\hline
6.0 & 1.7918 & 1.7934 & 1.7951 & 1.7967 & 1.7984 & 1.8001 & 1.8017 & 1.8034 & 1.8050 & 1.8066 \\
6.1 & 1.8083 & 1.8099 & 1.8116 & 1.8132 & 1.8148 & 1.8165 & 1.8181 & 1.8197 & 1.8213 & 1.8229 \\
6.2 & 1.8245 & 1.8262 & 1.8278 & 1.8294 & 1.8310 & 1.8326 & 1.8342 & 1.8358 & 1.8374 & 1.8390 \\
6.3 & 1.8405 & 1.8421 & 1.8437 & 1.8453 & 1.8469 & 1.8485 & 1.8500 & 1.8516 & 1.8532 & 1.8547 \\
6.4 & 1.8563 & 1.8579 & 1.8594 & 1.8610 & 1.8625 & 1.8641 & 1.8656 & 1.8672 & 1.8687 & 1.8703 \\
\hline
6.5 & 1.8718 & 1.8733 & 1.8749 & 1.8764 & 1.8779 & 1.8795 & 1.8810 & 1.8825 & 1.8840 & 1.8856 \\
6.6 & 1.8871 & 1.8886 & 1.8901 & 1.8916 & 1.8931 & 1.8946 & 1.8961 & 1.8976 & 1.8991 & 1.9006 \\
6.7 & 1.9021 & 1.9036 & 1.9051 & 1.9066 & 1.9081 & 1.9095 & 1.9110 & 1.9125 & 1.9140 & 1.9155 \\
6.8 & 1.9169 & 1.9184 & 1.9199 & 1.9213 & 1.9228 & 1.9242 & 1.9257 & 1.9272 & 1.9286 & 1.9301 \\
6.9 & 1.9315 & 1.9330 & 1.9344 & 1.9359 & 1.9373 & 1.9387 & 1.9402 & 1.9416 & 1.9430 & 1.9445 \\
\hline
7.0 & 1.9459 & 1.9473 & 1.9488 & 1.9502 & 1.9516 & 1.9530 & 1.9544 & 1.9559 & 1.9573 & 1.9587 \\
7.1 & 1.9601 & 1.9615 & 1.9629 & 1.9643 & 1.9657 & 1.9671 & 1.9685 & 1.9699 & 1.9713 & 1.9727 \\
7.2 & 1.9741 & 1.9755 & 1.9769 & 1.9782 & 1.9796 & 1.9810 & 1.9824 & 1.9838 & 1.9851 & 1.9865 \\
7.3 & 1.9879 & 1.9892 & 1.9906 & 1.9920 & 1.9933 & 1.9947 & 1.9961 & 1.9974 & 1.9988 & 2.0001 \\
7.4 & 2.0015 & 2.0028 & 2.0042 & 2.0055 & 2.0069 & 2.0082 & 2.0096 & 2.0109 & 2.0122 & 2.0136 \\
\hline
7.5 & 2.0149 & 2.0162 & 2.0176 & 2.0189 & 2.0202 & 2.0215 & 2.0229 & 2.0242 & 2.0255 & 2.0268 \\
7.6 & 2.0281 & 2.0295 & 2.0308 & 2.0321 & 2.0334 & 2.0347 & 2.0360 & 2.0373 & 2.0386 & 2.0399 \\
7.7 & 2.0412 & 2.0425 & 2.0438 & 2.0451 & 2.0464 & 2.0477 & 2.0490 & 2.0503 & 2.0516 & 2.0528 \\
7.8 & 2.0541 & 2.0554 & 2.0567 & 2.0580 & 2.0592 & 2.0605 & 2.0618 & 2.0631 & 2.0643 & 2.0656 \\
7.9 & 2.0669 & 2.0681 & 2.0694 & 2.0707 & 2.0719 & 2.0732 & 2.0744 & 2.0757 & 2.0769 & 2.0782 \\
\hline
8.0 & 2.0794 & 2.0807 & 2.0819 & 2.0832 & 2.0844 & 2.0857 & 2.0869 & 2.0882 & 2.0894 & 2.0906 \\
8.1 & 2.0919 & 2.0931 & 2.0943 & 2.0956 & 2.0968 & 2.0980 & 2.0992 & 2.1005 & 2.1017 & 2.1029 \\
8.2 & 2.1041 & 2.1054 & 2.1066 & 2.1078 & 2.1090 & 2.1102 & 2.1114 & 2.1126 & 2.1138 & 2.1150 \\
8.3 & 2.1163 & 2.1175 & 2.1187 & 2.1199 & 2.1211 & 2.1223 & 2.1235 & 2.1247 & 2.1258 & 2.1270 \\
8.4 & 2.1282 & 2.1294 & 2.1306 & 2.1318 & 2.1330 & 2.1342 & 2.1353 & 2.1365 & 2.1377 & 2.1389 \\
\hline
8.5 & 2.1401 & 2.1412 & 2.1424 & 2.1436 & 2.1448 & 2.1459 & 2.1471 & 2.1483 & 2.1494 & 2.1506 \\
8.6 & 2.1518 & 2.1529 & 2.1541 & 2.1552 & 2.1564 & 2.1576 & 2.1587 & 2.1599 & 2.1610 & 2.1622 \\
8.7 & 2.1633 & 2.1645 & 2.1656 & 2.1668 & 2.1679 & 2.1691 & 2.1702 & 2.1713 & 2.1725 & 2.1736 \\
8.8 & 2.1748 & 2.1759 & 2.1770 & 2.1782 & 2.1793 & 2.1804 & 2.1815 & 2.1827 & 2.1838 & 2.1849 \\
8.9 & 2.1861 & 2.1872 & 2.1883 & 2.1894 & 2.1905 & 2.1917 & 2.1928 & 2.1939 & 2.1950 & 2.1961 \\
\hline
9.0 & 2.1972 & 2.1983 & 2.1994 & 2.2006 & 2.2017 & 2.2028 & 2.2039 & 2.2050 & 2.2061 & 2.2072 \\
9.1 & 2.2083 & 2.2094 & 2.2105 & 2.2116 & 2.2127 & 2.2138 & 2.2148 & 2.2159 & 2.2170 & 2.2181 \\
9.2 & 2.2192 & 2.2203 & 2.2214 & 2.2225 & 2.2235 & 2.2246 & 2.2257 & 2.2268 & 2.2279 & 2.2289 \\
9.3 & 2.2300 & 2.2311 & 2.2322 & 2.2332 & 2.2343 & 2.2354 & 2.2364 & 2.2375 & 2.2386 & 2.2396 \\
9.4 & 2.2407 & 2.2418 & 2.2428 & 2.2439 & 2.2450 & 2.2460 & 2.2471 & 2.2481 & 2.2492 & 2.2502 \\
\hline
9.5 & 2.2513 & 2.2523 & 2.2534 & 2.2544 & 2.2555 & 2.2565 & 2.2576 & 2.2586 & 2.2597 & 2.2607 \\
9.6 & 2.2618 & 2.2628 & 2.2638 & 2.2649 & 2.2659 & 2.2670 & 2.2680 & 2.2690 & 2.2701 & 2.2711 \\
9.7 & 2.2721 & 2.2732 & 2.2742 & 2.2752 & 2.2762 & 2.2773 & 2.2783 & 2.2793 & 2.2803 & 2.2814 \\
9.8 & 2.2824 & 2.2834 & 2.2844 & 2.2854 & 2.2865 & 2.2875 & 2.2885 & 2.2895 & 2.2905 & 2.2915 \\
9.9 & 2.2925 & 2.2935 & 2.2946 & 2.2956 & 2.2966 & 2.2976 & 2.2986 & 2.2996 & 2.3006 & 2.3016 \\
\hline
$x$ & 0 & 0.01 & 0.02 & 0.03 & 0.04 & 0.05 & 0.06 & 0.07 & 0.08 & 0.09 \\ \hline
\end{tabular}\par}


\end{document}
